% =============================================================================
% ICLR 2027 submission -- negative results
% Plan of record: rebuild/PAPER/PAPER_PLAN.md
%
% Style files live in ../../iclr2027/ and are NEVER edited. Compile with:
%   cd rebuild/PAPER
%   TEXINPUTS=.:../../iclr2027: BIBINPUTS=.:../../iclr2027: pdflatex main.tex
%
% RULES INHERITED FROM THE REBUILD (see PAPER_PLAN.md):
%   * No number appears here unless a committed artifact supplies it. Every
%     factual sentence carries a source comment naming its log block / CSV.
%   * A number that does not exist is written [TODO: verify ...], never guessed.
%   * References are [CITE: topic] markers until a deep-research pass fills them.
%     Do NOT invent a citation key.
%
% ICLR 2027 SUBMISSION CHECKLIST (author to confirm before upload):
%   [ ] Main text ends on p.9 or earlier. Verified: Conclusion ends p.9. The AI
%       use / ethics / reproducibility statements, references and appendices are
%       excluded from the limit. Beyond 9 pages of main text = desk rejection.
%   [ ] \iclrfinalcopy stays COMMENTED OUT (anonymises authors, enables ruler).
%   [ ] No author identity in main text or supplementary material.
%   [ ] AI use statement completed and attested by an author. REQUIRED section.
%   [TODO: author] If any earlier version of this work exists on arXiv or was
%       presented at a workshop, cite it in the THIRD PERSON. Such a version does
%       not break anonymity, but referring to it as ours does.
%   [TODO: author] Supplementary material: measurement log, per-run metrics,
%       frozen decision rule, detector + self-test, contaminated filename list.
%       Anonymise any repository link; the host must not log visitors.
%   [ ] All co-authors have OpenReview profiles before the ABSTRACT deadline;
%       no authors may be added after it. Abstract must be genuine, not a
%       placeholder. At least one author registered to review >= 3 papers.
% =============================================================================
\documentclass{article}
\usepackage{iclr2027_conference,times}

% Optional math commands from https://github.com/goodfeli/dlbook_notation.
\input{math_commands.tex}

\usepackage{hyperref}
\usepackage{url}

% Added (not in the shipped template): figures and tables.
\usepackage{graphicx}
\usepackage{booktabs}
\usepackage{tikz}
\usetikzlibrary{positioning,arrows.meta,fit,backgrounds,calc}

%\iclrfinalcopy % Uncomment for camera-ready version, but NOT for submission.

\title{Uncertainty-Guided Closed-Loop Synthetic Generation\\
       Does Not Improve Synthetic-to-Real\\
       Camouflaged Object Detection}

% Overridden by the style file while \iclrfinalcopy is off; must be completed
% for the camera-ready version only.
\author{Anonymous}

\begin{document}

\maketitle

% =============================================================================
\begin{abstract}
% Every figure here is a back-reference; nothing originates in the abstract.
% 0.1600 -> Section 6 / Table 4; 41 of 76 and 53.9% -> Section 7 / Table 3;
% "within noise" -> Section 5 / Table 1. Reconciled 2026-09-09.
Camouflaged object detection is constrained by the cost of mask annotation,
which motivates training on synthetic imagery from camouflage generators. If
generation is cheap, a model's own uncertainty could steer \emph{where} the next
images are generated, closing a loop between training and data synthesis. We
test that idea. Instantiating the loop in a mean-teacher synthetic-to-real
framework with a published camouflage generator, we fix the decision rule before
training and run $24$ models across two architectures and three seeds. Targeted
allocation does not beat uniform allocation: the gap is within noise on both
architectures and both endpoints. We dissect the failure into four measured
causes and separate those specific to this framework and generator --- a
per-epoch optimisation budget that cannot grow with the dataset, and an
exhausted foreground supply whose objects are copied rather than synthesised ---
from those concerning the data and the signal, which travel further: the target
distribution's best silhouette coefficient is $0.1600$, and removing the
uncertainty information while preserving the allocation's shape reproduces the
same effect. A fifth candidate cause, too narrow a conditioning channel, is
refuted by source reading and withdrawn. Independently, we find that $41$ of
$76$ images in a standard COD
benchmark are re-encoded copies of training data, invisible to exact hashing,
and we release a detector and a clean evaluation protocol.
\end{abstract}

% =============================================================================
\section{Introduction}
\label{sec:intro}

Camouflaged object detection asks a model to segment objects whose appearance
has evolved to defeat segmentation. Ground-truth masks are correspondingly
expensive, and the field is small enough that most results rest on a handful of
datasets [CITE: COD10K] [CITE: NC4K]. Generative models offer an appealing way
out: a camouflage generator can synthesise a concealing background behind any
supplied foreground [CITE: LAKE-RED], and synthetic-to-real training then adapts
from generated data to real photographs [CITE: S2R-COD].

If generation is cheap, the natural next question is not \emph{how much} to
generate but \emph{where}. A trained model already knows which parts of the
target distribution it handles badly; that signal could steer the next batch of
generations toward exactly those regions, and the loop could be closed --- train,
measure deficiency, generate where deficient, retrain. The idea is well
motivated: it is the standard active-learning argument [CITE: uncertainty sampling] [CITE: generative active learning] applied to a generator instead of
an annotator, in a domain where annotation is the binding constraint. We
expected it to work.

\textbf{It does not.} We instantiate the loop in a mean-teacher
synthetic-to-real framework with a published camouflage generator, fix a decision
rule in advance, and train $24$ models across two architectures and three seeds.
Targeted allocation does not beat uniform allocation: the gap is within noise on
both architectures and both endpoints (Section~\ref{sec:null}).

A null is only worth reporting if it comes with a mechanism and an honest
account of its reach. We therefore separate our claims into three tiers and hold
to the distinction throughout. \textbf{T1} findings are properties of this
training loop or this generator; \textbf{T2} findings concern the target data or
the uncertainty signal and travel beyond our instantiation; \textbf{T3} findings
are independent of the method entirely. The distinction is not decoration: our
most striking cause --- that the training loop's optimisation budget cannot grow
with the dataset --- is T1, and reading it as evidence that synthetic data cannot
help would be unsupported by anything we measured.

\paragraph{Contributions.}
\begin{itemize}
\item \textbf{A pre-registered null.} Uncertainty-guided closed-loop generation
      yields no measurable accuracy gain in this setting. The decision rule was
      committed before the first training run and is reported unchanged,
      including the finding that the campaign resolved a coarser effect than it
      set out to (Section~\ref{sec:null}).
\item \textbf{A four-cause dissection, tiered by scope.} We localise the failure
      to a pinned optimisation budget (T1), an exhausted foreground supply
      (T1/T2), diffuse target structure (T2), and an uncertainty signal that
      adds little beyond concentration (T2); a fifth candidate, a conditioning
      bottleneck, is refuted by source reading and withdrawn
      (Section~\ref{sec:causes}).
% All figures in this list are back-references, not new measurements:
% 0.1600 -> Section 6 and Appendix Table 4 (EXP B1 #3);
% 41 of 76 / 53.9% -> Section 7 and Table 3 (EXP D2 #5, re-verified EXP D2R #4).
\item \textbf{A target distribution with little to target.} The target pool's
      best silhouette coefficient is $0.1600$ and stays below $0.25$ in all
      three embedding spaces we tried --- a property of the data, not of our
      trainer, and a precondition any cluster-allocation method needs and does
      not get here (Section~\ref{sec:causes}).
\item \textbf{Severe contamination in a standard COD benchmark, and a clean
      protocol.} Independently of the above, $41$ of $76$ CHAMELEON images
      ($53.9\%$) are re-encoded copies of training-pool images, invisible to
      exact hashing. We release a detector, the contaminated file list, and an
      evaluation protocol (Section~\ref{sec:contamination}).
\end{itemize}

% =============================================================================
\section{Related Work}
\label{sec:related}

\paragraph{Camouflaged object detection and its benchmarks.}
COD is evaluated on a small number of datasets --- COD10K, NC4K, CAMO and
CHAMELEON --- with architectures such as SINet, SINet-v2 and SegMaR
[CITE: COD10K] [CITE: NC4K] [CITE: CAMO] [CITE: CHAMELEON] [CITE: SINet]
[CITE: SINet-v2] [CITE: SegMaR]. Because these sets are reused across nearly all
published comparisons, their internal relationships matter for every reported
number; Section~\ref{sec:contamination} shows one such relationship has gone
unnoticed.

\paragraph{Synthetic data for camouflage.}
Annotating camouflaged objects is unusually expensive, which motivates
generating training data. Camouflage-specific generators paint a concealing
background behind a supplied foreground [CITE: LAKE-RED]
[CITE: camouflage generation other], typically built on diffusion inpainting
[CITE: diffusion inpainting]. Synthetic-to-real pipelines then adapt from
generated to real camouflage, of which the framework we test is a recent
instance [CITE: S2R-COD]. We test one generator inside one framework; other
generators are related work, not evidence about them.

\paragraph{Deciding what data to acquire.}
Choosing where to spend a labelling or generation budget is the subject of active
learning: uncertainty sampling, coverage and core-set criteria, and more recently
budget allocation over generative models
[CITE: active learning survey] [CITE: uncertainty sampling]
[CITE: core-set selection] [CITE: generative active learning]. Our stage (ii)
is an uncertainty criterion and stage (iii) an allocation rule, so the approach
under test is an active-learning policy over a generator rather than over an
annotator. Self-training with a mean teacher supplies the underlying adaptation
signal [CITE: mean teacher] [CITE: domain adaptation seg].

\paragraph{Negative results and benchmark integrity.}
Reporting methods that do not work is undersupplied relative to its value, and
pre-registration has been proposed to make such reports credible
[CITE: negative results venue] [CITE: pre-registration in ML]. Separately, a
line of work documents train--test contamination and near-duplicate leakage in
widely used datasets [CITE: train-test contamination]
[CITE: dataset duplication]. Our second contribution belongs to that line and,
unlike our first, does not depend on the method under test.

% =============================================================================
\section{The Approach Under Test}
\label{sec:method}

We specify the approach in full because our contribution is a precise account of
why it fails; each component below is later shown to be a distinct point of
breakdown. Section~\ref{sec:causes} maps each component to the measurement that
localises its contribution.

\subsection{Setting}

The base framework is cross-domain self-training with a mean-teacher objective
(CSRDA) for camouflaged object detection [CITE: S2R-COD]. A labelled
\emph{source} pool of synthetic camouflage images is paired with an unlabelled
\emph{target} pool of real photographs. A student network is trained by gradient
descent; a teacher is maintained as an exponential moving average of the student
with momentum $\lambda = 0.996$ and is never back-propagated.
% SOURCE: MyTrain.py:230 overwrites the argparse default 0.9996 with 0.996
% under --task S2C; update_ema at MyTrain.py:129. Verified this session.
Each optimisation step consumes one source batch and one target batch. The
objective is
\begin{equation}
\mathcal{L} \;=\; \mathcal{L}_{\mathrm{sup}} \;+\; \mathcal{L}_{\mathrm{ES}},
\qquad
\mathcal{L}_{\mathrm{ES}} \;=\; 0.9\,\mathcal{L}_{\mathrm{EA}} \;+\;
0.3\,\mathcal{L}_{\mathrm{SW}},
\label{eq:loss}
\end{equation}
where $\mathcal{L}_{\mathrm{sup}}$ is supervised loss on the labelled source
batch, $\mathcal{L}_{\mathrm{EA}}$ is a Sobel edge-alignment term between student
and teacher predictions on the target batch, and $\mathcal{L}_{\mathrm{SW}}$ is a
saliency-weighted cross-entropy against the teacher's prediction
[CITE: S2R-COD].
% SOURCE: coefficients a=0.9, b=0.3 set at MyTrain.py:225-231 under --task S2C;
% equation numbering (6-8) per Explanations/MEAN_TEACHER.md. Sobel kernels ->
% appendix.

Training runs for two self-training rounds. After the first round, a
confidence-based selection step (CLS) admits target images whose
student--teacher edge disagreement falls below a fraction of the running mean,
\begin{equation}
\mathrm{edge\_loss}(x) \;<\; u \cdot \overline{\mathrm{edge\_loss}},
\qquad u = 0.8,
\label{eq:cls}
\end{equation}
assigns each admitted image a pseudo-label by thresholding the teacher's class
activation map at $\tau = 0.4$, and appends the resulting pairs to the source
pool for the second round.
% SOURCE: CLS.py:139 (selection), CLS.py CAM threshold tau; u=0.8, tau=0.4 set
% at MyTrain.py:225-231. Round loop at MyTrain.py:265-348.
Model selection takes the teacher checkpoint minimising validation MAE on a
held-out real set, which is never used as an evaluation endpoint.
% SOURCE: MyTrain.py:181-188 (Tea_epoch_best rule), :302-304 / :322-324
% (--val_root ./Dataset/Val/CAMO/, 250 images asserted).

\subsection{The idea: closing the loop on the model's own uncertainty}

The synthetic source pool is produced by a background-generation model that
paints a novel camouflaging background behind a supplied foreground and its mask
[CITE: LAKE-RED]. Because generation is cheap relative to annotation, it is
natural to ask \emph{where} in the target distribution the next batch of
generations should be spent. The approach under test answers this with the
model's own error signal, in four stages.

\paragraph{(i) Target clustering.}
The unlabelled target pool is embedded with a frozen vision backbone and
partitioned by $k$-means into $k$ clusters. We use DINOv2-L at $518$px and
$k = 75$, selected by maximising the mean silhouette coefficient over
$k \in \{5,10,15,20,30,50,75,100,150\}$ [CITE: DINOv2] [CITE: silhouette].
% SOURCE: embedder identity = vit_large_patch14_dinov2.lvd142m at 518px
% (rebuild/E0/E0.md:142; rebuild/B1/out/b1_assertions.json family "dinov2").
% Exact model id -> appendix.
% SOURCE: rebuild/B1/out/b1_k_sweep.csv (silhouette curve, k*=75 at 0.1600);
% b1_embedder_sweep.json criterion = "max silhouette (stability reported, not
% used to rank)"; b1_cluster_assignment_dinoL518.json (k=75, seed=0).
% NOTE: the EXP B1 log NOTES prose describes a bootstrap-ARI criterion; the
% artifacts and B1_RESULTS.md:184 agree on silhouette. Artifacts are cited.
Embeddings are L2-normalised once at load time.
% SOURCE: rebuild/C1/c1_space.py load_space(); ABC_PLAN.md section A.3.

\paragraph{(ii) Deficiency scoring.}
Each cluster receives a scalar deficiency score, its mean per-image edge
disagreement $\mathit{es}$ under Eq.~\ref{eq:loss}, computed on the
\emph{target} images and requiring no ground truth. Scoring the target side
rather than a labelled endpoint is what makes the loop self-contained.
% SOURCE: rebuild/B1/out/b1_cluster_es_dinoL518.csv (75 rows, target_es
% populated 75/75, range 0.018391-0.070668) per ABC_PLAN.md section A.3.
% NOTE: EXP B1 block 4 established that the earlier blocks scored endpoint-side
% ES, which is not the signal the loop uses; allocation is by target_es.

\paragraph{(iii) Budget allocation.}
A generation budget of $B$ images is distributed over clusters by a softmax of
the deficiency scores at temperature $T = \alpha \cdot \mathrm{sd}(\mathit{es})$,
\begin{equation}
w_j \;\propto\; \exp\!\left(\frac{\mathit{es}_j - \max_i \mathit{es}_i}{T}\right),
\qquad
T = \alpha \cdot \mathrm{sd}(\mathit{es}),
\label{eq:alloc}
\end{equation}
converted to integers by the largest-remainder method so the allocation sums to
exactly $B$. We set $\alpha = 1.0$, the scale-free operating point at which the
temperature equals the spread of the scores themselves. At $B = 1000$ and
$\alpha = 1.0$ the allocation funds all $75$ clusters with a maximum share of
$0.194$.
% SOURCE: rebuild/C1/c1_targeted_vs_random.py softmax_alloc/largest_remainder,
% reused verbatim by ABC (ABC_PLAN.md section A.3, arm C); reproduction gate
% pins clusters_funded=75, max_alloc_share=0.194 against
% rebuild/C1/out/c1_cells.csv row (dinoL518, B=1000, R2_cut).

\paragraph{(iv) Foreground selection and generation.}
Within each funded cluster, foregrounds are served in order of cosine similarity
to the cluster centroid, and a greedy pass fills each cluster's quota, displacing
requests that cannot be met from the available foreground supply. The selected
foregrounds and their masks are passed to the generator, which returns one new
image per foreground. Selection is deterministic: it involves no random draw, so
the targeted condition is identical across training seeds.
% SOURCE: ABC_PLAN.md section A.3 (rank_by_centroid, serving_orders['desc_nc'],
% greedy_select); n_displaced=370 at alpha=1.0 per c1_cells.csv; stem list
% emitted once as rebuild/ABC/out/abc_stems_C_a1.0.txt and asserted identical
% across the three per-seed pools.

\subsection{Conditions}
\label{sec:conditions}

Isolating the effect of \emph{where} generations are spent requires holding
constant everything else that a larger pool changes. We therefore compare four
source pools that share the same base pool, verified byte-identical by per-file
SHA-256 against a provenance manifest rather than assumed:
% SOURCE: 1st EXP ABC block, threshold "every pool's base subset matches E0's
% manifest byte for byte" -> PASS; ABC_PLAN.md section A.5 item 7.

\begin{itemize}
\item \textbf{A0} --- base pool only. The condition comparable to the published
      baseline, and \emph{not} a clean control for the others: at a fixed
      optimisation budget a smaller pool gives every image more exposure per
      epoch (Section~\ref{sec:causes}).
\item \textbf{A2} --- base pool plus $B$ duplicates of the same foregrounds
      selected by condition B, carrying their \emph{original} backgrounds. This
      matches pool size and foreground identity while removing the generated
      background, isolating generation itself.
\item \textbf{B} --- base pool plus $B$ generated images for $B$ foregrounds
      drawn uniformly without replacement.
\item \textbf{C} --- base pool plus $B$ generated images for the $B$ foregrounds
      chosen by stages (i)--(iv). Identical injection mechanics to B; only the
      \emph{selection} differs.
\end{itemize}

\noindent
$\Delta(\mathrm{C}-\mathrm{B})$ is therefore the quantity of interest: targeted
against uniform allocation at equal pool size, equal budget, equal foreground
supply and equal mask provenance. $\Delta(\mathrm{B}-\mathrm{A2})$ separately
isolates the generated background.
% SOURCE: PREREGISTRATION.md section 1 defines Delta_CB as "THE claim".
% NOTE for later reconciliation: A2's added masks come from the original pool
% and B/C's from the generator's own output, a mask-tightness difference of
% 0.00575 mean white fraction on 1000 of 5447 images. This confounds
% Delta(B-A2) only; Delta(C-B) is unaffected (ABC_PLAN.md section A.3).

\subsection{What is held identical}

Beyond the base pool, the campaign fixes the architecture, the method flag, the
number of self-training rounds, the epoch and batch schedule, the learning rate
and decay, the loss coefficients of Eq.~\ref{eq:loss}, the seed governing
initialisation and shuffling, the unlabelled target pool, the model-selection
set and rule, the evaluation protocol, the generation budget $B$, and the
backbone initialisation, each asserted rather than assumed.
% SOURCE: ABC_PLAN.md section A.5 enumerates 15 held-identical items with the
% assertion used for each; the 1st EXP ABC block logs 16 thresholds, all PASS.
% Per-architecture schedule (epochs 40->39 / 100->99, batch 16/32, lr 1e-4,
% decay 0.1@30 / 0.1@50, clip_grad False/True) -> appendix, ABC_PLAN A.5.1.

One quantity is \emph{not} held constant by this design, and it is the reason the
null has a mechanism rather than merely an absence: the number of gradient steps
per epoch is set by the shorter of the two data loaders and is therefore pinned
by the target pool, independent of how large the source pool becomes. We return
to this in Section~\ref{sec:causes}.
% SOURCE: MyTrain.py:306-307 (min of loader lengths) with zip() truncation at
% :51,53; 2nd EXP ABC block threshold "total_step stayed pinned in BOTH rounds
% of every run (253 SINet / 127 SINet-v2)" -> PASS.

% =============================================================================
\section{Experimental Setup and Pre-Registration}
\label{sec:setup}

\paragraph{Endpoints.}
The primary endpoint is COD10K-test ($2026$ images), reported as the structure
measure $S_\alpha$ [CITE: S-measure] [CITE: COD10K]. NC4K ($4121$ images) is a
secondary endpoint: it is reported but never decides
[CITE: NC4K]. CHAMELEON is \emph{withdrawn} as an endpoint, for reasons
established independently of this campaign and reported in
Section~\ref{sec:contamination}. CAMO ($250$ images) is the model-selection set
and is never an endpoint.
% SOURCE: PREREGISTRATION.md section 1 (primary/secondary/never-decides);
% REBUILD_PLAN.md section 1 for counts; D2 scoping consequence "checkpoint
% selection runs on a test set => CAMO is never an endpoint".

\paragraph{Runs.}
Two architectures (SINet, SINet-v2) $\times$ four conditions $\times$ three
seeds $\{42,43,45\}$ give $24$ training runs, each of two self-training rounds.
% SOURCE: ABC_PLAN.md section A.2 (RUNID scheme); 2nd EXP ABC block
% "training_runs_completed = 24/24", "runs_abandoned = 0".
Because the campaign compares conditions rather than tuning them, the schedule
is fixed by the code and cannot drift from the command line; per-architecture
values are in Appendix~\ref{app:impl}.
% SOURCE: ABC_PLAN.md section A.5 item 4 and A.5.1; hyperparameters are
% overwritten after argument parsing at MyTrain.py:225-231.

\paragraph{The decision rule, fixed before the first run.}
The rule below was committed before any training began and never edited
afterwards; the governing document states that the campaign is void and must be
re-run if any part of it changed once a single $S_\alpha$ had been computed.
% SOURCE: rebuild/ABC/PREREGISTRATION.md, preamble and section 1. The 2nd and
% 3rd EXP ABC blocks both assert it unchanged.
Let $\hat{\sigma}$ be the pooled within-arm standard deviation of $S_\alpha$
across all four conditions and three seeds ($\mathrm{df} = 8$), estimated
separately for each architecture and endpoint. For a gap $\Delta$ between two
conditions:
\begin{center}
\begin{tabular}{ll}
\toprule
\textsc{real effect}     & $\Delta > 2\hat{\sigma}$ and sign consistent in $3/3$ seeds \\
\textsc{within noise}    & $|\Delta| \le 2\hat{\sigma}$ \\
\textsc{real regression} & $\Delta < -2\hat{\sigma}$ and sign consistent in $3/3$ seeds \\
\textsc{inconclusive}    & $|\Delta| > 2\hat{\sigma}$ but sign not $3/3$ --- reported as-is; seeds are \emph{not} added \\
\bottomrule
\end{tabular}
\end{center}
No $p$-value is computed at $n = 3$. The claim under test is
$\Delta(\mathrm{C}-\mathrm{B})$; $\Delta(\mathrm{B}-\mathrm{A2})$ is reported
alongside it, and gaps against A0 are reported for comparability with the
published baseline only.
% SOURCE: PREREGISTRATION.md section 1, quoted structure verbatim; the
% "do NOT add seeds" clause and "NO p-value at n=3" are in the frozen text.

\paragraph{Gates before any comparison.}
The pre-flight stage asserted $16$ conditions and all passed, including that
every condition's base pool matches the provenance manifest byte for byte, that
the targeted condition reproduces the allocation of the earlier selection
experiment exactly, that the two generated conditions overlap only at the chance
rate, and that the endpoint-measured signal appears nowhere in the pipeline.
% SOURCE: 1st EXP ABC block (2026-09-04T23:36:25+05:30), 16 THRESHOLD lines all
% -> PASS. NOTE: ABC_RESULTS.md reports "15/15" for this block; the block itself
% carries 16 THRESHOLD lines. The log wins (REBUILD_FINDINGS.md section 6).
The scoring code was itself validated against a previously committed
measurement before being used to decide anything.
% SOURCE: abc_verdict.json scorer_validation {Sm: 0.7172156230064066,
% MAE: 0.07446323197740462, passed: true} against
% Eval/Eval/eval_txt/SINet/S2C/10Aug_eval.txt.

% =============================================================================
\section{Result: No Measurable Gain}
\label{sec:null}

Table~\ref{tab:main} reports the campaign. \textbf{The claim under test,
$\Delta(\mathrm{C}-\mathrm{B})$, is \textsc{within noise} in all four
architecture $\times$ endpoint cells.} Targeting the generation budget at the
model's least-well-served clusters does not beat spending it uniformly, by the
rule fixed before the first run.
% SOURCE: 3rd EXP ABC block (2026-09-06T13:37:26+05:30); abc_verdict.json
% verdicts.<cell>.gaps["B->C10"].verdict == "WITHIN NOISE" in all four cells.

Two features of the table deserve comment before any interpretation. First, the
arm means \emph{do} increase monotonically from A0 through A2 and B to C in
three of the four cells, and the ordering is the one the approach predicts --- it
is simply smaller than the noise the design was built to see through. Second,
the largest single gap anywhere in the campaign is not the one under test.
% SOURCE: arm_means in abc_verdict.json; monotonicity noted in
% A3_RESULTS.md section 6 ("the arm means in fact rose monotonically in all four
% architecture x endpoint cells").

\begin{table}[t]
\caption{\textbf{Main result.} $S_\alpha$ per condition, mean $\pm$ standard
deviation over seeds $\{42,43,45\}$; higher is better. $\Delta(\mathrm{C}-\mathrm{B})$
is the pre-registered claim and $2\hat{\sigma}$ its decision threshold; ``sign''
in parentheses counts seeds agreeing in direction. All values at four decimal places. The frozen
decision rule, the per-run metrics and the machine-readable verdict are in the supplementary
material; Appendix~\ref{app:trace} maps every value to its measurement-log entry.}
\label{tab:main}
\begin{center}
\footnotesize
\begin{tabular}{llccccrc}
\toprule
Arch. & Endpoint & A0 & A2 & B & C & $\Delta(\mathrm{C}-\mathrm{B})$ & $2\hat{\sigma}$ \\
\midrule
SINet    & COD10K$^\dagger$ & 0.7010 & 0.7077 & 0.7134 & 0.7185 & $+0.0050$ (3/3) & 0.0179 \\
SINet    & NC4K             & 0.7581 & 0.7632 & 0.7669 & 0.7690 & $+0.0021$ (3/3) & 0.0083 \\
SINet-v2 & COD10K$^\dagger$ & 0.6891 & 0.6921 & 0.6951 & 0.6947 & $-0.0004$ (2/3) & 0.0123 \\
SINet-v2 & NC4K             & 0.7444 & 0.7479 & 0.7486 & 0.7496 & $+0.0009$ (2/3) & 0.0105 \\
\midrule
\multicolumn{8}{c}{\textsc{verdict: within noise in all four cells}} \\
\midrule
\multicolumn{8}{l}{Per-arm s.d., primary cell: A0 $0.0173$, A2 $0.0019$, B $0.0018$, C $0.0041$.} \\
\multicolumn{8}{l}{$^\dagger$ primary endpoint. $|\Delta| \le 2\hat{\sigma}$ in every row.} \\
\bottomrule
\end{tabular}
\end{center}
\end{table}

\paragraph{The one significant gap is inadmissible, and we say so rather than report it.}
Across all $20$ gap tests the campaign returns \textsc{within noise} eighteen
times, \textsc{inconclusive} once, and \textsc{real effect} once. The single
\textsc{real effect} is $\mathrm{A0} \rightarrow \mathrm{C}$ on SINet/NC4K,
$\Delta = +0.0109$, consistent in $3/3$ seeds. It is disqualified twice over: it
sits on the secondary endpoint, which the pre-registered rule states never
decides, and it is measured against A0, which is not a clean control because at a
pinned optimisation budget its smaller pool gives every image more exposure per
epoch (Section~\ref{sec:causes}). We report it for completeness and draw nothing
from it.
% SOURCE: abc_verdict.json -- verdict census 18 WITHIN NOISE / 1 INCONCLUSIVE /
% 1 REAL EFFECT; the REAL EFFECT is SINet|NC4K gaps["A0->C10"], delta 0.010892,
% n_seeds 3. ABC_RESULTS.md:71-74 warns these "must not be read as support".
% REAL REGRESSION never occurs.

\paragraph{Power: the design did not reach the resolution it declared.}
The pre-registration stated that the design ``resolves an effect $\sim$half the
paper's MT$\rightarrow$Ours gap ($0.0142$) and no smaller''. The realised noise
estimate on the primary endpoint is $2\hat{\sigma} = 0.0179$, which
\emph{exceeds} that reference gap: the campaign resolves roughly $1.3\times$ it,
not half of it. \textbf{The null is therefore a null at a coarser resolution
than intended, and both halves of that sentence belong in any report of this
result.}
% SOURCE: ABC_RESULTS.md:6-11 and :120-123 ("it resolves about 1.3x it").
% 2*sigma_hat = 0.017933 at REBUILD_LOG.txt:1453. The 0.0142 reference is
% Sa 0.7172 (Eval/.../SINet/S2C/10Aug_eval.txt) - 0.7030
% (Eval/.../SINet/S2C_MT/12Aug_eval.txt).
The rule anticipated this: it fixes the bar at $2\hat{\sigma}$ as measured, not
at an assumed noise level, so a larger $\hat{\sigma}$ raises the bar rather than
invalidating the test.
% SOURCE: PREREGISTRATION.md section 2.7 pre-authorises a rising bar.
Two quantities that would sharpen the statement are unavailable and are not
estimated here: run-to-run seed variance, which would require retraining that
was deferred, and any rate translating a pool shift into an $S_\alpha$ change,
which we found could not be re-derived from the available runs and therefore do
not assume.
% SOURCE: sigma_seed is UNVERIFIED-DEFERRED (REBUILD_PLAN.md C3, REVISION_TABLE
% section 5). The old response-rate anchor 0.00867 Sa/SD is recorded as
% discarded and non-derivable in REBUILD_PLAN.md section 3 (C3): "I will not
% invent a rate."

% =============================================================================
\section{Why It Fails}
\label{sec:causes}

A null is only useful if it comes with a mechanism. We measured four, and they
do not have the same reach. A fifth, a conditioning bottleneck in the generator,
we advanced and withdraw in paragraph~(ii). Recalling the scope tiers of
Section~\ref{sec:intro}: \textbf{T1} findings are properties of this training
loop or this generator, \textbf{T2} findings concern the data or the signal and
travel beyond our instantiation, and \textbf{T3} is independent of the method
entirely. Table~\ref{tab:causes} summarises; the tier column is the load-bearing
one.

\begin{table}[t]
\caption{\textbf{Four measured causes, with scope, and one refuted.} T1 is
specific to this framework or generator; T2 concerns the data or signal and
generalises further. A T1 cause is \emph{not} evidence about
generators-in-the-loop in general. Row (ii), refuted by source reading, is
withdrawn and carries no tier. Roman numerals index
the paragraphs of Section~\ref{sec:causes};
Appendix~\ref{app:trace} gives the measurement-log entry for each row.}
\label{tab:causes}
\begin{center}
\footnotesize
\begin{tabular}{@{}p{0.265\linewidth}p{0.545\linewidth}c@{}}
\toprule
Cause & Measured quantity & Tier \\
\midrule
(i) Optimisation budget is pinned & Steps per epoch fixed at $253$ / $127$ in both rounds of all $24$ runs, independent of pool size & T1 \\
(ii) Generator conditioning width \emph{--- refuted, withdrawn} & Not a bottleneck: the foreground latent also reaches the denoiser at full resolution, where global self-attention can read it from any background position & --- \\
(iii) Foreground pool is exhausted & Generated set is a bijection onto the $4447$ existing foregrounds; object pixels are copied, not synthesised & T1/T2 \\
(iv) Target clusters are diffuse & Best silhouette $0.1600$; below $0.25$ in all three embedding spaces & T2 \\
(v) Signal adds little over shape & Removing the targeting information but keeping the allocation shape reproduces the same separation & T2 \\
\bottomrule
\end{tabular}
\end{center}
\end{table}

\paragraph{(i) Added data cannot buy optimisation. \textnormal{[T1]}}
Each step consumes one source and one target batch, and the loop stops when the
shorter loader is exhausted. Because the target pool is fixed, the number of
gradient steps per epoch is pinned --- measured at $253$ (SINet) and $127$
(SINet-v2) in \emph{both} rounds of all $24$ runs, regardless of how much
synthetic data was added.
% SOURCE: 2nd EXP ABC block, THRESHOLD "total_step stayed pinned in BOTH rounds
% of every run (253 SINet / 127 SINet-v2)" -> PASS; abc_runs.csv total_step_set.
The consequence is \emph{dilution, not exclusion}: both loaders shuffle, so the
added images are sampled --- they occupy about $18.4\%$ of the first round's
gradient content --- but every image's expected exposure per epoch falls
correspondingly. A condition with the base pool alone gives each image $0.9103$
exposures per epoch against $0.7432$ for the enlarged pools, some $22\%$ more.
% SOURCE: 2nd EXP ABC block NOTES ("0.910 ... against 0.743 ... 22% more");
% POOL_MECHANICS_AUDIT.md section 7(a) for the dilution-not-exclusion
% correction and the 18.4% round-1 share.
% NOTE: that audit has no EXP block of its own; its round-1 figures are also in
% the log block, and its round-2 figures are NOT logged -> not used here.
This is the cause most easily over-read. It says that \emph{this} loop converts
extra data into a rearranged mixture rather than more learning; it says nothing
about whether generated data helps a training regime that scales its
optimisation with dataset size. It also explains why A0 is not a clean control,
and hence why A2 exists.

\paragraph{(ii) The conditioning channel --- a cause we withdraw. \textnormal{[T1, refuted by source reading]}}
A selection method can only matter if the difference it selects for survives
into the render, and we advanced a narrow conditioning channel as the reason it
might not: before the foreground conditions background synthesis it is reduced
to $16$ superpixel means of a three-channel latent --- $48$ scalars --- which
would bound what any selection policy could express. Reading the generator's
source refutes that account. The summary is real but is one route among several.
The conditioning tensor handed to the diffusion denoiser retains the foreground
latent at full resolution over the object region and enters as input channels of
a U-Net whose self-attention is unmasked and global, so any background position
can read the foreground directly; a second route admits the same
full-resolution latent into every regenerated position through a $1\times1$
convolution inside the summarising module. Foreground information is therefore
not confined to the summary, and we withdraw the cause.
% SOURCE: static reading of the generator; the exact files read are hashed in
% rebuild/A1/out/a1_source_manifest.sha256 and the audit is rebuild/A1/A1_SCOPING.md.
% 48 = n_super_pix 16 (config_LAKERED.yaml:72) x embed_dim 3 (:47).
% Summary construction: ddpm.py:1548-1564 (LMP), :1574-1586 (BKRM + upsample).
% Route via the 1x1 conv: ddpm.py:1534 (Conv2d(6,3,kernel_size=1)) and :1587
% (fuse over cat(vec_bg, fg)). Route via the U-Net: ddpm.py:1590 (new_fg keeps
% fg over the object region) -> :1467 (concatenated to the U-Net input;
% in_channels 7, config:29) -> openaimodel.py:318-324 (AttentionBlock flattens
% all spatial positions, no mask), instantiated at :541 and :606 for
% attention_resolutions [8,4,2] (config:32-35).
% NOTE: this paragraph reports no measured quantity. It is a reading of source,
% not a measurement, and it therefore has no measurement-log entry; the manifest
% above is what pins it, because LAKE-RED/ is untracked in this repository.
% SUPERSEDES the earlier "[T1, unmeasured]" wording, which named the decisive
% check as whether the foreground tensor is identically zero under the mask.
% That check is both unsatisfiable -- the first-stage encoder normalises over the
% whole spatial extent (GroupNorm, model.py:38-39), so every latent position is a
% function of global input statistics -- and insufficient, since the U-Net route
% above does not pass through the tensor it inspects.
It does not follow that the generator makes \emph{semantic} use of what it
receives --- that needs a foreground perturbation with the summary held fixed,
which we did not run --- nor that it is steerable
(Section~\ref{sec:conclusion}). The withdrawal removes a T1 explanation and
leaves more of the account to (iv) and (v).
% SOURCE: absence of a cross-attention port -- openaimodel.py:464-475 default
% use_spatial_transformer=False / context_dim=None, and unet_config
% (config_LAKERED.yaml:25-43) sets neither; the concat branch calls the denoiser
% with no context argument (ddpm.py:1467-1468). The only conditioning is the
% 4-channel masked-image latent plus mask (test.py:128-134).
% NOTE: "does not establish semantic use" is the A1 limit recorded in
% rebuild/A1/A1_SCOPING.md section 4.3; the perturbation design is its section 4.2.

\paragraph{(iii) There are no new foregrounds to find. \textnormal{[T1 mechanism, T2 consequence]}}
Every image the loop can add is a re-rendering of a foreground already in the
pool. The generated set is a bijection onto the $4447$ source foregrounds:
$4447$ of $4447$ trace back, none originates outside, and none is left
unrendered. Moreover the generator composites the original object pixels back
into its output, so the object is reproduced rather than re-imagined --- mean
absolute error inside the object is $5.603$ against $70.802$ in the background
(a ratio of $12.64$), and not one of the $4447$ images has an object-interior
error above $40$.
% SOURCE: EXP D1 block; D1_RESULTS.md sections 3.1-3.2. The second pool
% (authors') gives 9.633 vs 70.555, ratio 7.32; both reported in the appendix.
Exhaustion here is therefore an identity statement about pixels, not a
statistical claim about distributional overlap that could be argued about. The
mechanism is specific to a generator that composites the object back; the
consequence --- that re-rendering an exhausted foreground supply adds no new
object information --- applies to any pipeline whose novelty is confined to
backgrounds.

\paragraph{(iv) The target distribution has little cluster structure to target. \textnormal{[T2]}}
The method assumes the target pool decomposes into clusters that can be
individually under-served. It barely does. The best mean silhouette coefficient
over $k \in \{5,\dots,150\}$ is $0.1600$, at $k = 75$ --- far below the
$0.25$ that would indicate usable structure --- and the finding is robust across
embedding spaces: $0.1465$ for DINOv2-L at $224$px and $0.0568$ for CLIP-L, with
the peak below $0.25$ in every space.
% SOURCE: EXP B1 block 3, THRESHOLD "WEAK CLUSTER STRUCTURE is embedder-robust:
% the silhouette peak is below 0.25 in every embedder space" -> PASS;
% b1_k_sweep*.csv, b1_embedder_sweep.json. Full grid in Appendix Table 4.
We report a companion failure honestly: in the CLIP space the silhouette does
not attain an interior maximum at all but falls monotonically to $0.0357$ at
$k = 150$, so ``choose $k$ by silhouette'' is not even well posed there.
% SOURCE: EXP B1 block 3, THRESHOLD "the silhouette criterion actually selects
% a k in every embedder space (an interior maximum, not the grid edge)" -> FAIL.
This is a property of the target data, not of our trainer: an allocation policy
over clusters is only as meaningful as the clustering it allocates over.

\paragraph{(v) The uncertainty signal contributes little beyond concentration. \textnormal{[T2]}}
The targeted and uniform conditions do produce measurably different pools. But
an audit of \emph{why} they differ found that the difference is not attributable
to the uncertainty signal: destroying the targeting information while preserving
the \emph{shape} of the allocation --- how concentrated it is, not where it
points --- reproduces the same separation. Of seven declared attribution
thresholds, two passed and five failed, and the honest verdict is that the
separation is real but is produced by concentration rather than by targeting.
% SOURCE: 4th EXP C1 block; C1_RESULTS.md section 8 ("the ES signal is not what
% produces d ~ 1.0"), verdict REOPENS-BUT-NOT-BY-TARGETING, 2/7 PASS 5 FAIL.
This matters more than the null itself: it means the negative result is not
merely about a weak generator, but about a signal that, in this setting, carries
little information a concentration prior does not already carry.

\paragraph{Where the synthetic images sit, and a metric that cannot tell us. \textnormal{[T2]}}
It is tempting to explain the null by saying the synthetic pool is simply far
from the real target distribution. Stated absolutely, that explanation does not
survive its own controls. A linear probe separates real from generated images at
an area-under-curve of $0.9993$ --- but re-encoding the \emph{identical} real
images as JPEG at quality $30$ separates them from themselves at $0.9928$,
higher than either genuine real-versus-real control ($0.9381$ and $0.9225$).
The declared vacuity flag fired: at that operating point the metric is not
measuring distributional distance.
% SOURCE: EXP A3 block, THRESHOLD "T6 VACUITY FLAG: some identity-preserving
% control exceeds AUC 0.90" -> PASS (it fired); A3_RESULTS.md:118-170.
What does survive is comparative and stated as coverage: a genuinely different
real camouflage dataset covers the target manifold at $0.79$--$0.90$, real
photographs of an unrelated genre at $0.71$--$0.75$, and the synthetic images
generated \emph{from those very photographs} at $0.13$--$0.54$.
% SOURCE: A3_RESULTS.md:280-284, quoted as the strongest defensible form.
We state the limit of this explicitly: it is a characterisation of the
distribution, not a bound on training utility. A distant or low-coverage
synthetic set can still improve a model, and in our own campaign the condition
means rose monotonically. A3 supplies a mechanism consistent with the absence of
a resolvable gain; it does not show that no gain exists.
% SOURCE: A3_RESULTS.md section 6 "What A3 does NOT establish", verbatim intent.

% =============================================================================
\section{Benchmark Contamination and a Clean Evaluation Protocol}
\label{sec:contamination}

This section is independent of everything above. It would stand if the method
under test had worked, and it is why CHAMELEON is absent from
Table~\ref{tab:main}. \textbf{Of the $76$ images in CHAMELEON, $41$
($53.9\%$) are re-encoded copies of images in the training pool a
COD10K-trained model has already seen} --- $40$ in COD10K-train and one in CAMO.
% SOURCE: 41/76 (53.9%) first logged in EXP D2 block 3 and carried to the
% authoritative block 5; re-measured against an author-sourced copy in EXP D2R
% block 4 (41/76, OLD CLAIM -> MATCH, delta 0). The 40/1 partner split is D2R
% only (nearvdup partner_pool).

\paragraph{Exact hashing declares the set clean.}
Hashing CHAMELEON against the training pool returns \emph{zero} collisions at
the file-byte level and zero at the decoded-pixel level. The duplicates are
re-encodings: same photograph, different compression. Any audit that stops at
hashing --- as audits usually do --- will certify a $53.9\%$-contaminated
benchmark as uncontaminated. That is the trap this section exists to close.
% SOURCE: EXP D2 block 5, "cham vs target: 0 byte, 0 pixel"; the log line notes
% "0 here does not mean clean".

\paragraph{Detection.}
Images are grouped by exact pixel dimensions, since a re-encode preserves them.
Within a group, each image is reduced to a $32\times32$ greyscale descriptor
and pairs are shortlisted by descriptor RMS distance; survivors are decoded at
full resolution and confirmed when the mean absolute difference over RGB falls
below a tolerance of $6/255$. We deliberately do \emph{not} contrast-normalise
the descriptor: normalising destroys the scale that discriminates a re-encode
from a merely similar photograph, and it is precisely what caused our own first
pass to report $10$ of $76$ rather than $41$.
% SOURCE: detect_contamination.py:64-69 (constants), :206-216 and :281-298
% (grouping, shortlist, confirmation), :17-20 (the rejected normalisation).
% The 10/76 (13.2%) figure is EXP D2 blocks 1-2; superseded at block 3.
% Complexity: 23,854 images -> 4,427 dimension groups -> 7.4M candidate pairs
% -> 323 shortlisted (d2_leakage_sweep.py:198-202; D2_RESULTS.md:205-212).
The criterion is a \emph{mean}, not an exact match: the confirmed pairs reach
maximum per-pixel differences as high as $70$, so any test demanding pixel
equality rejects all of them.
% SOURCE: d2r_duplicate_pairs.csv max_abs column, range 13-70.

\paragraph{The count does not depend on the tolerance.}
Sweeping the confirmation tolerance over $\{1,2,3,5,6\}$ gives
$11, 26, 37, 40, 41$ contaminated images: the count saturates rather than
growing with permissiveness. Independently of any tolerance, the sorted
nearest-neighbour distances show a gap --- $41$ images lie below $5.51$ and the
next is at $40.58$, a factor of $7.36$ --- so the contaminated set is separated
from the rest of the benchmark by an order of magnitude, not by a threshold we
chose. Varying the shortlist depth changes nothing.
% SOURCE: EXP D2 block 5 and EXP D2R block 4 tolerance sweeps (identical for
% both copies); NN gap below=5.512 above=40.577 n_below=41 ratio 7.36x;
% TOPK sweep {8,32,all} all give 41 (d2r_topk_sensitivity.json).
% The other three endpoints show no gap.

\begin{table}[t]
\caption{\textbf{Contamination against the COD10K training pool.} ``Exact''
counts byte- or decoded-pixel-identical files, ``re-enc.'' counts confirmed
re-encodings at tolerance $6/255$; the two are disjoint by construction. The
nearest-neighbour gap is tolerance-independent. Lower panel: the CHAMELEON count
as the confirmation tolerance varies. Appendix~\ref{app:trace} gives the measurement-log entry
for each row.}
\label{tab:contam}
\begin{center}
\footnotesize
\begin{tabular}{@{}lrrrrl@{}}
\toprule
Benchmark & $n$ & Exact & Re-enc. & Total (\%) & NN gap \\
\midrule
CHAMELEON                  & 76   & 0 & 41 & \textbf{41 (53.9)} & $41$ below $5.51$, next $40.58$ ($7.36\times$) \\
COD10K-test                & 2026 & 7 & 2  & 9 (0.44)           & none --- continuous \\
NC4K                       & 4121 & 0 & 1  & 1 (0.02)           & none --- continuous \\
CAMO$^\ddagger$            & 250  & 0 & 4  & 4 (1.6)            & none --- continuous \\
\midrule
NC4K vs COD10K-\emph{test} & 4121 & 0 & 0  & \textbf{0 (0.0)}   & none; closest $20.61$, $3.44\times$ tol. \\
\midrule
\multicolumn{6}{@{}l}{CHAMELEON count at tolerance $1/2/3/5/6$: \; $11 \,/\, 26 \,/\, 37 \,/\, 40 \,/\, \mathbf{41}$ \; (saturates)} \\
\midrule
\multicolumn{6}{@{}l}{$^\ddagger$ model-selection set, never an endpoint. Rows other than CHAMELEON are measured} \\
\multicolumn{6}{@{}l}{against on-disk copies, not author-sourced ones; treat them as lower bounds.} \\
\bottomrule
\end{tabular}
\end{center}
\end{table}

\paragraph{The finding survives being re-tested on data we did not control.}
Our first audit used the CHAMELEON copy in our own working tree. Because that is
exactly the weakness such a claim invites, we repeated it against the release
distributed by the dataset's authors. The two copies are byte-identical in all
$76$ files, and the re-audit returns the same $41$ of $76$. It also corrected
three of our own earlier values --- we report those corrections rather than
quietly adopting them.
% SOURCE: EXP D2R block 4; copies byte-identical 76/76; OLD CLAIM -> MATCH for
% 41/76; three MISMATCH entries recorded as corrections to our own prior work.

\paragraph{One declared check failed, and we leave it failing.}
We had declared that every confirmed pair should show differing JPEG
quantisation tables as independent evidence of re-encoding. It does not hold:
two files in the CHAMELEON release are PNGs carrying a \texttt{.jpg} extension
and so have no quantisation table at all, and our original comparison had read a
\emph{missing} table as a differing one --- absence of evidence counted as
evidence. The corrected tally is $40$ of $41$ pairs evidenced by quantisation
table and one by container format ($413$\,KB PNG against a $132$\,KB JPEG), so
every pair retains independent evidence, but the check as worded fails and we
did not reword it after the fact.
% SOURCE: EXP D2R block 4, THRESHOLD T5 -> FAIL in all four D2R blocks; the
% replacement T5b -> PASS ("every confirmed pair carries re-encoding evidence
% from SOME independent channel"). Files: animal-19.jpg, animal-28.jpg.
% D2R_RESULTS.md:110-120.

\paragraph{What we do \emph{not} claim: inflated scores.}
The obvious next step is to argue that reported CHAMELEON numbers are inflated.
We measured this and cannot support it. Removing the seven exact duplicates from
COD10K-test moves MAE by $1.2 \times 10^{-5}$, and in the wrong direction; the
duplicated images sit at the $48$th percentile of the clean set's error
distribution. On CHAMELEON, the contaminated and clean subsets are
indistinguishable on every proxy we tried, and the direction of the difference
\emph{flips} depending on which of the two circulating mask releases is used to
score. \textbf{The conclusion is about independence, not inflation: a
contaminated column cannot be interpreted as an independent measurement,
regardless of whether it happens to be higher.}
% SOURCE: EXP D2 block 5 MAE impact (delta -1.242e-05, -0.0167%, mean
% percentile 0.4761); EXP D2R block 4 leaked-vs-clean (all percentiles
% 0.448-0.559, Mann-Whitney p 0.47-0.99,
% split_direction_agrees_across_mask_sets = False). CLEAN_PROTOCOL.md:145-148
% states the independence-not-inflation conclusion; the D2R inference run is
% explicitly not quotable as a performance number.

\paragraph{A second, larger source of incomparability.}
While reconciling the two CHAMELEON mask releases we found that the choice of
release alone changes MAE by a factor of $2.7$ on \emph{identical} predictions
($0.2196$ against $0.0816$), with only $27$ of $76$ masks identical between them.
Papers reporting CHAMELEON without naming the mask release are not comparable
with each other even setting contamination aside.
% SOURCE: EXP D2R block 4, d2r_reconcile.json; mean IoU 0.6932 after polarity
% alignment, 0.0 as stored, opposite polarity conventions.

\paragraph{A clean protocol.}
For a COD10K-trained camouflaged-object model we recommend: report COD10K-test
and NC4K, disclosing that both were audited against on-disk copies and that a
fraction of each is unchecked; do not report CHAMELEON as an independent
endpoint, and if it is reported at all, report it on the named $35$-image
uncontaminated subset together with the mask release used; never report the set
used for model selection as an endpoint; and run the detector on your own
directories before reporting anything. NC4K in particular survives the
additional check its use invites --- it does not overlap COD10K-test at any
level, with the closest approach $3.44\times$ the confirmation tolerance.
% SOURCE: CLEAN_PROTOCOL.md:202-217 (recommendation, quoted in substance) and
% :32-37 (the reportability table); EXP D2_NC4K block 2 for the NC4K result.
% NOTE: cite the repository CLEAN_PROTOCOL.md, not release/CLEAN_PROTOCOL.md,
% which omits the EXP D2_NC4K section entirely.
The detector, the list of contaminated filenames, and the protocol are included in the
supplementary material.
% [TODO: author -- if an anonymous repository is preferred to a supplementary
% .zip, insert the anonymised link here. It must not reveal author identity and
% the host must not log visitors (ICLR 2027 author guidelines).]
Two caveats are ours to carry: unchecked is not clean, so every rate in
Table~\ref{tab:contam} is a lower bound; and only the CHAMELEON row has been
verified against author-sourced data, so the others remain provisional.
% SOURCE: CLEAN_PROTOCOL.md:45-54 (both universal caveats, verbatim intent).
% NOTE: our own protocol table lists COD10K-test as 2/2026; that is the
% re-encode count only and omits the 7 exact duplicates, which are disjoint.
% Table 3 above reports the corrected total of 9/2026 (0.44%).
% [TODO: fix CLEAN_PROTOCOL.md:34 upstream to 9/2026.]

% =============================================================================
\section{Conclusion}
\label{sec:conclusion}

Uncertainty-guided closed-loop generation does not improve synthetic-to-real
camouflaged object detection in the setting we tested. The gap between targeted
and uniform allocation is within noise on two architectures and two endpoints
under a rule fixed before training, and the campaign resolved a coarser effect
than it was designed to. We hold the scope of that result where the measurements
put it: the two causes with the largest apparent force --- a per-epoch
optimisation budget that cannot grow with the dataset, and a foreground supply
whose objects are copied rather than synthesised --- are properties of this loop
and this generator. What generalises further is less flattering to the idea
itself: the target distribution has little cluster structure to allocate over,
and the uncertainty signal contributes little beyond the concentration of the
allocation it induces.

\paragraph{What a working system would require.}
The causes imply concrete requirements rather than general prohibitions.
\emph{First}, a training regime whose optimisation scales with dataset size, so
that added data buys gradient steps instead of diluting exposure; under a fixed
step budget no acquisition policy can express itself. \emph{Second}, a source of
genuinely new foregrounds. Our null is correctly read as ``targeting does not
help once the foreground pool is exhausted'', and the experiment that would
separate exhaustion from targeting --- a condition supplied with new foregrounds
from a salient-object dataset, holding the generator fixed --- is designed and
deferred rather than run, because admitting a pool that has not been swept for
leakage to the standard of Section~\ref{sec:contamination} into a load-bearing
comparison would repeat the error that section documents.
% SOURCE: ABC_SCOPING.md section 8 (the Oracle arm: DECISION "OUT of this
% campaign ... blocked on data that is not on this machine"; its scope is a
% clean test of foreground exhaustion holding the generator fixed; the four
% unblocking prerequisites include a D2-standard leakage sweep). D1_RESULTS.md
% section 5: "D1 does not establish that new foregrounds could not help."
\emph{Third}, a deficiency signal aligned with what the task actually rewards.
Ours is an edge-disagreement measure, and its own attribution audit found the
targeting information largely redundant given the allocation shape.
\emph{Fourth}, a generator that can be steered toward a named deficiency rather
than only supplied with a foreground --- a requirement about the \emph{shape} of
the interface, not its capacity. Ours carries foreground information at full
resolution (Section~\ref{sec:causes}(ii)) but accepts only that foreground and
its mask, with no cross-attention port through which a deficiency could be
named. Whether a generator offering one would help is untested.

\paragraph{What survives regardless.}
The contamination result does not depend on any of the above. A benchmark in
routine use is $53.9\%$ re-encoded training data, and exact hashing certifies it
as clean. The detector, the contaminated file list, and the evaluation protocol
are released, and we expect them to outlast the negative result they were found
alongside.

% =============================================================================
\subsection*{AI use statement}

(This section is \textbf{required} and does not count toward the page limit.)

% [TODO: author must read, correct and attest to this section before submission.
% It is drafted from what actually occurred; only an author can confirm the
% review claims in the final two sentences. Required section; do not delete.]
We used a large language model in three roles, and disclose each.
\emph{Analysis.} An LLM was used to query, cross-check and summarise the
measurement artifacts described in Appendix~\ref{app:trace} --- for example, to
recompute per-condition means from the committed per-run metrics and to check
reported figures against the measurement log. It did not perform any of the
measurements themselves: every number in this paper was produced by a committed
script and written to the log before being quoted.
\emph{Drafting and editing.} An LLM assisted in drafting and revising the
manuscript text, including the structure of Sections~\ref{sec:method}--\ref{sec:conclusion}
and the tables, working from the committed artifacts and a written plan.
\emph{Not used.} We did not use generative AI to produce research ideas, to
generate or synthesise any data or experimental result, or to write the
experiment code whose outputs this paper reports; the training, audit and
detection code predates and is independent of the drafting process.
We reviewed all AI-assisted work: every quantitative claim was checked against
its source artifact, and claims we could not source are marked as unverified in
Appendix~\ref{app:ledger} rather than asserted. We take responsibility for the
final content of this work, including all text, claims and artifacts.

\subsection*{Ethics statement}

This work reports that a widely used evaluation benchmark substantially overlaps
a standard training split, and releases the means to detect it. We note three
considerations. \emph{First}, the finding concerns dataset construction, not
conduct: the overlap is invisible to exact hashing, so a benchmark can be
assembled and used in good faith without anyone noticing, and we attribute no
fault to the dataset's authors or to any group that has reported results on it.
\emph{Second}, we deliberately do not convert the finding into a claim about any
published number. We measured whether contamination inflates reported scores and
found no detectable effect (Section~\ref{sec:contamination}); asserting that
specific results are inflated would go beyond our evidence and could unfairly
impugn other work. \emph{Third}, the released material --- a detector and a list
of filenames already present in a public dataset --- introduces no new personal
data and involves no human subjects. We release it because the alternative,
reporting a contamination rate without the means to reproduce or extend the
check, would ask the community to take the claim on trust. The method we test is
a third party's; we test one instantiation of it, scope our conclusions
accordingly, and report the negative result without implying error on the part
of its authors.

\subsection*{Reproducibility statement}
The experimental setup, the pre-registered decision rule and the gates asserted
before any comparison are given in Section~\ref{sec:setup}; per-architecture
schedules, backbone and embedder identifiers, and the exposure calculation are
in Appendix~\ref{app:impl}. Appendix~\ref{app:trace} maps every quantity
reported in the paper to the entry in the measurement log that produced it, and
Appendix~\ref{app:ledger} states the threshold accounting and lists every
quantity we could not verify. The detection procedure of
Section~\ref{sec:contamination} is described in enough detail to reimplement,
and the implementation, its self-test and the affected filenames are provided in
the supplementary material together with the measurement log, the per-run
metrics and the frozen decision rule.

\bibliography{iclr2027_conference}
\bibliographystyle{iclr2027_conference}

\appendix

\section{Implementation details}
\label{app:impl}

\paragraph{Schedules.}
Hyperparameters are set in code after argument parsing and cannot be overridden
from the command line, so they are identical across all conditions. SINet: $40$
epochs declared and $39$ executed, batch $16$, learning rate $10^{-4}$ decayed
by $0.1$ at epoch $30$, no gradient clipping. SINet-v2: $100$ declared and $99$
executed, batch $32$, learning rate $10^{-4}$ decayed by $0.1$ at epoch $50$,
gradient clipping enabled.
% SOURCE: ABC_PLAN.md section A.5.1, restated from MyTrain.py:194-199, :249-259.

\paragraph{Steps and exposure.}
Steps per epoch are $253$ (SINet) and $127$ (SINet-v2), giving $4048$ and $4064$
source images seen per epoch. Per-image expected exposure per epoch is $0.9103$
for the base-only condition against $0.7432$ for the enlarged pools on SINet, and
$0.9139$ against $0.7461$ on SINet-v2.
% SOURCE: ABC_PLAN.md section A.5.1; 2nd EXP ABC block; abc_runs.csv.

\paragraph{Backbones and embedders.}
Backbone initialisations are pinned by file size and SHA-256
(\texttt{resnet50-11ad3fa6.pth}, $102{,}540{,}417$ bytes;
\texttt{res2net50\_v1b\_26w\_4s-3cf99910.pth}, $103{,}197{,}949$ bytes).
Embeddings use \texttt{vit\_large\_patch14\_dinov2.lvd142m} at $224$ and $518$px
and \texttt{vit\_large\_patch14\_clip\_224.openai}, all frozen, with features
L2-normalised at use.
% SOURCE: ABC_PLAN.md section A.5 item 15; rebuild/E0/E0.md:141-142;
% rebuild/B1/out/b1_assertions.json.

\paragraph{Object-compositing measurement, both pools.}
Mean absolute error inside the object against the background was $5.603$ vs
$70.802$ (ratio $12.64$) for the locally regenerated pool and $9.633$ vs
$70.555$ (ratio $7.32$) for the pool the trainer actually reads, over all $4447$
images of each, with zero objects exceeding an interior error of $40$.
% SOURCE: EXP D1 block; D1_RESULTS.md section 3.2.

\paragraph{Deferred by design.}
Resize interpolation, the numerical stabiliser in the allocation softmax,
tie-breaking in the greedy selection, and the Sobel kernels of
Eq.~\ref{eq:loss} are unchanged from the released implementations of the
framework and generator and are not restated here.
% [TODO: verify -- readable from the released code, but not pinned by a
% committed measurement in this repository.]

\section{Cluster quality across $k$ and embedding space}
\label{app:clusters}

\begin{table}[h]
\caption{\textbf{Mean silhouette coefficient} over $k$-means partitions of the
$4040$-image target pool, by embedding space and $k$. The peak of each row is in
bold. No space exceeds $0.25$; in CLIP the curve has no interior maximum,
falling across the grid, so selecting $k$ by silhouette is not well posed there.
Full sweep artifacts are in the supplementary material; see Appendix~\ref{app:trace}.}
\label{tab:silhouette}
\begin{center}
\footnotesize
\setlength{\tabcolsep}{3.5pt}
\begin{tabular}{@{}lccccccccc@{}}
\toprule
Embedding space & $k{=}5$ & $10$ & $15$ & $20$ & $30$ & $50$ & $75$ & $100$ & $150$ \\
\midrule
DINOv2-L / 518px & 0.0566 & 0.0923 & 0.1152 & 0.1326 & 0.1527 & 0.1556 & \textbf{0.1600} & 0.1557 & 0.1507 \\
DINOv2-L / 224px & 0.0538 & 0.0868 & 0.1053 & 0.1180 & 0.1378 & \textbf{0.1465} & 0.1457 & 0.1419 & 0.1367 \\
CLIP-L / 224px   & \textbf{0.0568} & 0.0550 & 0.0486 & 0.0486 & 0.0488 & 0.0448 & 0.0396 & 0.0382 & 0.0357 \\
\bottomrule
\end{tabular}
\end{center}
\end{table}

\section{Traceability}
\label{app:trace}

Every quantity in this paper originates in an append-only measurement log
supplied with the supplementary material. One entry is written per accepted run
and records the command, the environment, the declared thresholds with their
outcomes, and the artifacts produced; entries are never edited or deleted, so
superseded runs remain visible alongside the ones reported. Table~\ref{tab:trace}
maps each reported group of quantities to its entry.

\begin{table}[h]
\caption{\textbf{Where each reported quantity comes from.} Entry labels are those
used in the supplied measurement log. Where several entries share a label, the
one cited is named; earlier entries are retained but superseded. One row reports
a reading of the generator's source rather than a measurement and so has no log
entry; the files read are hashed in the supplementary material.}
\label{tab:trace}
\begin{center}
\footnotesize
\begin{tabular}{@{}p{0.42\linewidth}p{0.30\linewidth}p{0.22\linewidth}@{}}
\toprule
Reported in & Quantity & Log entry \\
\midrule
Section~\ref{sec:setup}, gates & $16$ pre-flight assertions & \texttt{ABC} \#1 \\
Table~\ref{tab:main}; Section~\ref{sec:null} & condition means, $\hat{\sigma}$, gaps, verdicts & \texttt{ABC} \#3 \\
Section~\ref{sec:causes}(i); App.~\ref{app:impl} & steps per epoch, exposure & \texttt{ABC} \#2 \\
Section~\ref{sec:causes}(ii) & conditioning-path architecture & source reading; no entry \\
Section~\ref{sec:causes}(iii); App.~\ref{app:impl} & foreground bijection, compositing error & \texttt{D1} \\
Section~\ref{sec:causes}(iv); Table~\ref{tab:silhouette} & silhouette across $k$ and space & \texttt{B1} \#3 \\
Section~\ref{sec:causes}(v) & attribution audit & \texttt{C1} \#4 \\
Section~\ref{sec:causes}, last paragraph & coverage, probe controls, vacuity flag & \texttt{A3} \\
Table~\ref{tab:contam}, CHAMELEON row & contamination rate, tolerance sweep & \texttt{D2} \#5; re-audit \texttt{D2R} \#4 \\
Table~\ref{tab:contam}, NC4K vs COD10K-test & independence of the two endpoints & \texttt{D2\_NC4K} \#2 \\
Section~\ref{sec:contamination}, mask releases & MAE under each mask release & \texttt{D2R} \#4 \\
App.~\ref{app:ledger} & threshold accounting & all entries cited above \\
\bottomrule
\end{tabular}
\end{center}
\end{table}

\section{Measurement discipline and what remains unverified}
\label{app:ledger}

Across the log blocks cited in this paper, $143$ thresholds were declared in
advance, $127$ passed and $16$ failed. Failures are reported rather than
repaired: several are the findings themselves (the contamination rate, the
instability of the binary cluster-structure label), two are corrections to our
own earlier work, and one --- the quantisation-table check of
Section~\ref{sec:contamination} --- was left failing rather than reworded after
the fact.
% SOURCE: REBUILD_FINDINGS.md sections 1-3 (cited-block ledger 143/127/16).

\noindent The following are \emph{not} established by any committed measurement.
Except where a reading of source is named as the basis, they are asserted
nowhere in this paper:
\begin{itemize}
\item Whether the generator makes semantic use of the foreground information it
      receives. That the architecture does \emph{not} confine that information
      to the $48$-scalar summary we establish by reading the generator's source
      rather than by measurement (Section~\ref{sec:causes}(ii)); the experiment
      that would show whether the available capacity is \emph{used} ---
      perturbing the foreground while holding the summary values fixed and
      measuring the generated background --- was not run. We assert no bound in
      either direction.
\item Any inflation of reported scores attributable to contamination. We
      measured this and found no detectable effect; the conclusion is about
      independence, not inflation.
\item Run-to-run seed variance, and any rate translating a pool shift into a
      change in $S_\alpha$. Both would require training that was deferred; the
      latter could not be re-derived from the available runs.
\item Author-sourced verification of the contamination rates for COD10K-test,
      NC4K and CAMO. Only the CHAMELEON row has been re-audited against data we
      did not control. [TODO: verify --- re-audit owed.]
\item Whether any published table reports the contaminated benchmark, which we
      cannot check from the artifacts available to us.
\item An overlap we measured but did not log to the standard the rest of this
      work requires: a further $32$ of the $76$ CHAMELEON images appear to
      coincide with COD10K-\emph{test}. [TODO: verify --- present in the
      artifact, absent from the measurement log; used in no claim above.]
\end{itemize}

\end{document}
